\section{Conclusion}
\label{sec:conclusion}

We revisited the distinction between within-event interpolation and
event-held-out generalization in SAR water segmentation. The corrected
comparison uses identical official test chips for each surviving
event-excluded checkpoint and a fixed generalist. On nine recoverable
event--seed pairs across seven events, five differences favored the
event-excluded checkpoint and four favored the generalist. Dropping
one chip could change the sign in seven pairs. This dependence on the
observed chips, together with incomplete seeds and different training
and validation sets, limits what the comparison says about new events.

The Azerbaijan prototype uses SegFormer-B4 and applies JRC occurrence
and terrain/coverage exclusions to the raw water masks. The two
historical scene products show how strongly those choices affect the
mapped extent, particularly where the DEM is missing. Accuracy in
Azerbaijan remains unmeasured. Repeated training on matched designs
and independent inundation labels at the satellite overpass would
address the benchmark and prototype limitations, respectively.

\paragraph{Code and reproducibility.}
The historical benchmark uses ML commit \texttt{867a4f3}. The upload
API and offline CLI use the
\href{https://github.com/AliAgabalayev/Flood-Damage-Detection/commit/23563a7bebf28e508b394cf6858649718fd148a3}{terrain-corrected ML commit}
(\texttt{23563a7}). The
\href{https://github.com/AliAgabalayev/flood-paper/tree/paper-v1.0}{companion manuscript repository}
provides archived per-chip confusion counts, scene replay records and
commands to regenerate the tables and figures.
